% Generated draft: requires mathematical and source review.
\documentclass[11pt]{article}
\usepackage[margin=25mm]{geometry}
\usepackage{fontspec}
\setmainfont{Latin Modern Roman}
\usepackage{amsmath,amssymb,mathtools,unicode-math}
\setmathfont{Latin Modern Math}
\usepackage{enumitem,needspace,xurl,hyperref}
\setlength{\parindent}{0pt}
\setlength{\parskip}{5pt}
\setlength{\emergencystretch}{3em}
\newcommand{\N}{\mathbb{N}}
\newcommand{\Z}{\mathbb{Z}}
\newcommand{\Q}{\mathbb{Q}}
\newcommand{\R}{\mathbb{R}}
\newcommand{\C}{\mathbb{C}}
\newcommand{\F}{\mathbb{F}}
\newcommand{\A}{\mathbb{A}}
\newcommand{\PP}{\mathbb{P}}
\newcommand{\E}{\mathbb{E}}
\newcommand{\eps}{\varepsilon}
\newcommand{\dd}{\,\mathrm{d}}
\begin{document}
\section*{UnsolvedMath ID: 30004620}
\textbf{Source problem code:} catalogue entry
\subsection{Definitions and mathematical statement}
Let \(\mathcal A_3\) denote the moduli space of principally polarized abelian varieties of dimension three, and let \(X\) denote the toroidal compactification intended by the source. The supplied record does not specify the base field, the toroidal decomposition, or whether a stack or its coarse moduli space is intended.

Under a numerical-cycle interpretation, let \(N_2(X)_{\R}\) be the real vector space of algebraic cycles of dimension two modulo numerical equivalence, and define
\[
C_2(X)=\left\{\sum_{i=1}^{r}a_i[V_i]:\ r\in\Z_{\geq0},\ a_i\in\R_{\geq0},\ V_i\subseteq X\text{ an integral closed surface}\right\}.
\]
Here \([V_i]\) denotes the numerical class of \(V_i\), and the empty sum is zero. The supplied record does not specify whether its notation \(\operatorname{Eff}_2\) means \(C_2(X)\), its closure \(\overline{C_2(X)}\), or a cone in another cycle-class space.

Write \(s_A,s_F,s_D,s_{K3+1},s_{C4}\) for the classes intended by the source symbols \(S_A,S_F,S_D,\sigma_{K3+1},\sigma_{C4}\), respectively. Their geometric definitions are not included in the supplied record. With these unresolved choices retained, the surface-cone question in the supplied clean statement is
\[
\operatorname{Eff}_2(X)\stackrel{?}{=}
\R_{\geq0}s_A+
\R_{\geq0}s_F+
\R_{\geq0}s_D+
\R_{\geq0}s_{K3+1}+
\R_{\geq0}s_{C4}.
\]
Equivalently, under the numerical-cycle interpretation and the source's assertion that the five classes are effective, the question asks whether every class \(\alpha\in\operatorname{Eff}_2(X)\) admits coefficients \(a_A,a_F,a_D,a_{K3+1},a_{C4}\in\R_{\geq0}\) such that
\[
\alpha=a_As_A+a_Fs_F+a_Ds_D+a_{K3+1}s_{K3+1}+a_{C4}s_{C4}.
\]
The record writes \(\operatorname{Eff}_2(\mathcal A_3)\) while describing cycles on a compactification; \(X\) makes that intended ambient space explicit without selecting a compactification.

There is a material conflict in the supplied formulations: the original statement asks about a codimension-two effective cone, whereas the clean statement asks about surfaces. In the usual genus-three moduli setting, \(\dim\mathcal A_3=6\); surfaces have codimension four, and codimension-two cycles have dimension four. Thus these formulations are not interchangeable. The displayed question records the supplied clean surface formulation, but the intended target, the five geometric classes, and the cone convention require confirmation before this becomes a fully specified mathematical statement.
\subsection{Short English statement}
The supplied clean formulation asks whether every effective surface class on the intended toroidal compactification of \(\mathcal A_3\) is a nonnegative real combination of five named classes. The compactification, class definitions, and cone convention are missing, and the original codimension-two wording conflicts with the surface formulation. The supplied open status and extremality assessment are dated curation, not independently verified conclusions.
\subsection{Sources}
\begin{itemize}
\item Local UnsolvedMath registry, \texttt{data/unsolvedmath/problems.json}, permanent ID \(30004620\), exact source problem code \texttt{catalogue entry}, title: Effective Surface Cone of Principally Polarized Abelian Threefolds. The supplied record belongs to the Oberwolfach Reports Open Problems collection, gives proposed year \(2021\), and names no proposer.
\item Supplied source citation: \emph{Moduli spaces and Modular forms (2022)}, Oberwolfach Reports, pp.~301--359, DOI \texttt{10.4171/owr/2021/6}. \url{https://doi.org/10.4171/owr/2021/6}. The background describes Oberwolfach Reports as documenting mathematical workshops and their open-problem discussions. The supplied citation year is retained without resolving its relationship to the year in the DOI or the proposed year.
\item Supplied literature source: Grushevsky--Hulek, effective surfaces. \url{https://arxiv.org/abs/2007.02995}. This is the reference label supplied by the record; no additional bibliographic details are supplied here.
\item Supplied literature source: Oberwolfach report. \url{https://doi.org/10.4171/owr/2021/6}.
\item Literature assessment dated exactly \texttt{2026-08-22}: the supplied curation labels the problem open, states that the five displayed effective rays are known to be extremal, and says that equality with the effective surface cone remains conjectural. This is supplied dated curation derived from the OWR material, not an independently verified literature assessment. The source explicitly requests independent verification of time-sensitive status and cited claims.
\item The registry labels the statement \texttt{corrected\_verified} and reports that it was checked against the official Oberwolfach report. That is supplied verification metadata, not verification performed for this draft, and it does not resolve the formulation issues described above.
\end{itemize}
\end{document}
